\documentclass[12pt]{article}
\usepackage[utf8]{inputenc}
\usepackage[T1]{fontenc}
\usepackage{hyperref}
\title{Kuamasi Language Resources for K-12 Curriculum Development}
\author{Compiled by Claude Sonnet 5.5 (claude-sonnet-5-5) for Andr\'as Kornai}
\date{}
\begin{document}
\maketitle

\section{Introduction}
Kuamasi (ISO 639-3: yku; Glottocode kuam1234) is a Loloish (Sino-Tibetan) language of Yunnan, China. It is also called Songping Yi, and its Chinese name is the one given by Wikipedia. It is spoken in Matang and Songping Village, Liuhe Township, Heqing County. Wikipedia gives about 1,000 native speakers (2011), and Joshua Project lists about 1,100 people. The search found no evidence that it is taught in schools.

\section{Primary Educational Resources}
The search found no primers, textbooks or school materials.

\section{Linguistic and Reference Materials}
\begin{itemize}
\item Wikipedia article on the Kuamasi language: \url{https://en.wikipedia.org/wiki/Kuamasi_language}
\item Glosbe language page for Kuamasi (yku). Its content and quality were not verified: \url{https://glosbe.com/yku/nmx}
\end{itemize}

\section{Religious and Community Resources}
\begin{itemize}
\item Joshua Project language page for yku. It reports no Bible translation and an unknown writing status: \url{https://joshuaproject.net/languages/yku}
\end{itemize}

\section{Digital and Audio Resources}
The search found no audio, video or software resources.

\section{Curriculum Development Context}
Almost nothing is documented, and nothing pedagogical was found. Any curriculum would first need community-led documentation and a practical orthography, and the community's own wishes about the language should come first.

\end{document}